\documentclass[10pt]{article}
\usepackage[margin=0.8in]{geometry}
\usepackage[hidelinks]{hyperref}
\usepackage{enumitem}
\setlist{nosep, leftmargin=*}
\setlength{\parskip}{3pt}
\setlength{\parindent}{0pt}

\begin{document}

{\large\bfseries DS-US 202 Project Proposal: Technical Audit of the Home Credit Default Risk ADS}\\[2pt]
\textbf{Team:} [Member 1 (NetID)], [Member 2 (NetID)], [Member 3 (NetID)], [Member 4 (NetID)] \hfill October 14, 2026

\paragraph{ADS under audit.}
We will audit a credit default risk model built for the \emph{Home Credit Default Risk} Kaggle competition~\cite{competition}.
Home Credit is a consumer lender that targets applicants with little or no credit history. The ADS takes a loan application
and outputs a score for the risk that the applicant will have payment difficulties (\texttt{TARGET}=1). The lender uses this
score to decide whether to approve the loan. The competition ranked submissions by ROC AUC on a hidden test set.

\paragraph{Data.}
We use the official competition data~\cite{competition}, downloaded from Kaggle under the competition rules. The audited
implementation uses the main application table: \texttt{application\_train.csv} (307,511 applications $\times$ 122 columns,
8.07\% positive) and \texttt{application\_test.csv} (48,744 applications, unlabeled). The features include demographics
(sex, age, family status, education, number of children), income and loan amounts, employment and housing information, and
three normalized scores from external data sources (\texttt{EXT\_SOURCE\_1--3}). The competition also provides auxiliary tables
(credit bureau records, previous applications, installment histories) that the baseline does not use.

\paragraph{Implementation.}
We audit the public notebook \emph{Start Here: A Gentle Introduction} by Will Koehrsen (version~17)~\cite{koehrsen}.
\begin{itemize}
  \item \textbf{Preprocessing:} label and one-hot encoding of categorical features; a flag for the anomalous
        \texttt{DAYS\_EMPLOYED}~=~365243 code; median imputation; min--max scaling; optional polynomial and domain-knowledge features.
  \item \textbf{Models:} a logistic regression, random forests, and a LightGBM gradient-boosted model~\cite{lightgbm}
        trained with 5-fold cross-validation and class weighting. We will focus our audit on the LightGBM model
        with domain-knowledge features, the final model in the notebook.
  \item \textbf{Reproduction:} we ran the entire notebook locally from start to finish without errors.
        We changed only deprecated library calls for current
        scikit-learn and LightGBM versions; the modelling logic is unchanged.
\end{itemize}

\paragraph{Why this ADS.}
Credit approval is a high-stakes decision. This model uses attributes that anti-discrimination law protects in lending (such as
sex, age and marital status under the U.S.\ Equal Credit Opportunity Act) directly as input features, which lets us audit both
group fairness and disparate treatment. It also relies heavily on opaque third-party scores, which raises the transparency and
interpretability questions studied in this course.

\paragraph{Preliminary audit plan.}
\begin{itemize}
  \item \textbf{Subgroups:} sex, age bands, family status, and education, plus their intersections.
  \item \textbf{Accuracy by group:} AUC, TPR/FPR at decision thresholds, and calibration. The model's mean predicted score
        (0.35) far exceeds the base rate (8\%), so its scores are not calibrated probabilities and we must interpret them with care.
  \item \textbf{Fairness:} demographic parity, equalized odds, and predictive parity, computed with Fairlearn~\cite{fairlearn}.
  \item \textbf{Robustness and interpretability:} stability across folds and random seeds, feature attributions,
        and ablations that remove protected attributes and the external scores.
\end{itemize}

\vspace{-6pt}
\begin{thebibliography}{9}\small
\setlength{\itemsep}{0pt}
\bibitem{competition} A.~Montoya, inversion, K.~Odintsov, M.~Kotek. \emph{Home Credit Default Risk}. Kaggle, 2018.
  \url{https://www.kaggle.com/competitions/home-credit-default-risk}
\bibitem{koehrsen} W.~Koehrsen. \emph{Start Here: A Gentle Introduction} (version 17). Kaggle notebook, 2018.
  \url{https://www.kaggle.com/code/willkoehrsen/start-here-a-gentle-introduction}
\bibitem{lightgbm} G.~Ke et al. LightGBM: A Highly Efficient Gradient Boosting Decision Tree. \emph{NeurIPS}, 2017.
\bibitem{fairlearn} H.~Weerts et al. Fairlearn: Assessing and Improving Fairness of AI Systems. \emph{JMLR} 24(257), 2023.
\end{thebibliography}

{\small\textbf{AI Disclosure:} We selected a suitable project from Kaggle's beginner resources and manually ran the code locally in
VS Code. We then used Claude to review our outputs and help check the project's feasibility and whether the model's
hyperparameter settings were suitable for the baseline analysis.\par}

\end{document}
